\documentclass[10pt]{article}

\usepackage[utf8]{inputenc}

\usepackage[T1]{fontenc}

\usepackage{amsmath}

\usepackage{amsfonts}

\usepackage{amssymb}

\usepackage[version=4]{mhchem}

\usepackage{stmaryrd}



\begin{document}

хов, 1941, получил его из модельного уравнения для спектра турбулентности).



Форму функции $\varphi_{1}(\lambda k)$ при больших $k\left(\gg \lambda^{-1}\right)$ можно определить, учитывая, что поле скорости в масштабах $l \ll \lambda$ приблизительно линейно и в системе координат $x_{\beta}$, движущейся и вращающейся вместе с жидкой частицей и связанной с ее главными осями деформации, имеет вид $u_{\beta}=a_{\beta} x_{\beta}$, где $0<a_{1} \geqslant a_{2} \geqslant a_{3}<0$ - главные скорости деформации. Тогда решение линеаризованного уравнения для вихря получается в виде



\[

\omega=A(t) \exp [i x k(t)]

\]



причем при больших $\boldsymbol{t}$ вектор $\boldsymbol{k}$ становится параллельным оси $x_{3}$ и имеющим длину



\[

k_{3}(0) \exp \left(\left|a_{3}\right| t\right)

\]



a вектор $\omega$ - параллельным оси $x_{1}$ и имеющим длину



Отсюда спектр энергии



\[

\left|A_{1}(0)\right| \exp \left(a_{1} t-v \int_{0}^{t} k^{2} d t\right)

\]



получается в виде



\[

E(k)=\left\langle\delta \omega \delta \omega^{*} / 2 k^{2} \delta k\right\rangle

\]



\[

E(k) \sim \varepsilon^{2 / 3} k^{-5 / 3}(\lambda k)^{2 \sigma-4 / 3} \exp \left[-(\gamma \lambda k)^{2}\right],

\]



где $\sigma=\left\langle a_{1}\left|a_{3}\right|^{-1}\right\rangle$ и $\gamma^{2}=\tau^{-1}\left\langle\left|a_{3}\right|^{-1}\right\rangle(\tau-$ колмогоровский внутренний масштаб времени), причем $1 / 2<\sigma<1$, и можно положить $\sigma=2 / 3$.



Турбулентности ускорения мелкомасштабны (определяются в основном движениями масштабов $l \approx \lambda$ ), имеют средний квадрат порядка $\varepsilon^{3 / 2} v^{-1 / 2}$ и при движении жидкой частицы в инерционном интервале времен практически некоррелированы (А. М. Яглом, 1949). Поле вихря скорости также мелкомасштабно, имеет средний квадрат $\varepsilon / v$ и спектр в инерционном интервале $C_{1} \varepsilon^{2 / 3} k^{1 / 3}$.



Концентрации $\theta$ пассивных примесей (в том числе, если архимедовы силы несущественны, то также и температуры) в поле локально изотропной турбулентности при достаточно болыших числах Пекле $\mathrm{Pe}=\chi^{-1} L_{\theta} \delta_{\theta} U$, где $\chi$ - кинематич. коэффициент диффузии, $L_{\theta}$ - типичный пространственный масштаб существенных изменений значений $\boldsymbol{\theta}$, $\delta_{\theta} U-$ типичный перепад скоростей на расстояниях $L_{\theta}$, допустимо считать локально изотропными. Для них первая гипотеза подобия должна быть заменена (А. М. Обухов, 1949) предположением, что конечномерные распределения вероятностей для разностей скоростей и разностей концентраций полностью определяются четырьмя размерными параметрами v, $\varepsilon, \chi и \varepsilon_{\theta}$ (автомодельность по Re и $\mathrm{P} \mathrm{e})$, где $\varepsilon_{\theta}-$ скорость переноса удельной меры $(1 / 2)\left\langle\theta^{\prime 2}\right\rangle$ неоднородностей $\theta^{\prime}=\theta-\langle\theta\rangle$ поля концентрации по спектру масштабов от больших масштабов к малым (равная удельной скорости $\chi\left\langle\left|\nabla \theta^{\prime}\right|^{2}\right\rangle$ выравнивания этих неоднородностей вследствие молекулярной диффузии). При этом возникает внутренний масштаб поля концентрации



\[

\lambda_{\theta}=\left(\chi^{3} / \varepsilon\right)^{1 / 4}=\lambda(\mathrm{Pr})^{-3 / 4},

\]



где $\operatorname{Pr}=v / \chi$ - число Прандтля. Согласно сформулированной гипотезе пространственная структурная функщия поля $\theta$ при $r \ll L, L_{\theta}$ и спектр этого поля при $k \gg L^{-1}, L_{\theta}^{-1}$ должны иметь вид



\[

\begin{aligned}

& \left\langle(\delta \theta)^{2}\right\rangle=\varepsilon_{\theta} \varepsilon^{-1 / 3} r^{2 / 3} \varphi_{\theta}\left(r / \lambda_{\theta}, \operatorname{Pr}\right) \\

& E_{\theta}(k)=\varepsilon_{\theta} \varepsilon^{-1 / 3} k^{-5 / 3} \varphi_{1 \theta}\left(\lambda_{\theta} k, \operatorname{Pr}\right),

\end{aligned}

\]



где $\varphi_{\theta}\left(x_{\theta}, \operatorname{Pr}\right)$ и $\varphi_{1 \theta}\left(x_{\theta}, \operatorname{Pr}\right)$ - нек-рые универсальные функции. При очень болыших $\mathrm{Pe}$ появляется конвективный интервал масштабов $L, L_{\theta} \gg r \gg \lambda_{\theta}^{*}$, в к-ром сглаживающее влияние диффузии на поле концентрации мало по сравнению с ростом градиентов этого поля, создаваемым перено- сом примеси движущейся жидкостью (конвекцией), так что параметром $\chi$ в этом интервале можно пренебречь. Оказывается, что при $\operatorname{Pr} \ll 1$ конвективный интервал простирается только до масштабов



\[

l \gg \lambda_{\theta}

\]



а при $\operatorname{Pr} \gg 1$ - до масштабов



\[

l \gg\left(v \chi^{2} / \varepsilon\right)^{1 / 4}=\left(\lambda \lambda_{\theta}^{2}\right)^{1 / 3} .

\]



На инерционно-конвективном интервале



\[

\min \left(L, L_{\theta}\right) \gg r \gg \max \left(\lambda, \lambda_{\theta}\right)

\]



можно пренебречь обоими параметрами $\chi$ и v, и тогда



\[

\begin{gathered}

\left\langle(\delta \theta)^{2}\right\rangle=C_{\theta} \varepsilon_{\theta} \varepsilon^{-1 / 3} r^{2 / 3}, \\

E_{\theta}(k)=C_{1 \theta} \varepsilon_{\theta} \varepsilon^{-1 / 3} k^{-5 / 3}

\end{gathered}

\]



Первая из этих формул принадлежит А. М. Обухову (1949), вторая - закон пяти третей для температурного поля - С. Корpсину (S. Corrsin, 1951). По эмпирич. данным, $C_{1 \theta} \approx 1,1$.



В случае $\operatorname{Pr} \gg 1$ при больших $k\left(\gg \lambda^{-1}\right)$ форму спектра $E_{\theta}(k)$ можно установить при помощи гипотезы Бэтчелора, по к-рой здесь этот спектр может зависеть от $\varepsilon$ лишь через типичную величину $\tau^{-1}=(\varepsilon / v)^{1 / 2}$ скоростей деформации. В вязкоконвективном интервале



\[

1 \ll \lambda k \ll(\operatorname{Pr})^{1 / 2}

\]



коэффициент диффузии $\chi$ влиять на вид спектра еще не должен, и получается промежуточная асимптотика



\[

E_{\theta}(k) \sim \varepsilon_{\theta} \tau k^{-1} .

\]



Лишь затем в вязкодиффузионном интервале



\[

\lambda k \gg(\mathrm{Pr})^{1 / 2}

\]



из-за диффузионного сглаживания спектр резко спадает по закону



\[

k^{-1} \exp \left[-(\gamma \lambda k)^{2} / \operatorname{Pr}\right]

\]



В случае же $\operatorname{Pr} \ll 1$ в интервале масштабов



\[

\lambda_{\theta} \gg l \gg \lambda

\]



поле концентрации приблизительно линейно $\left(\theta \approx b_{\alpha} x_{\alpha}\right)$, и уравнение диффузии достаточно брать в виде



\[

b_{\alpha} u_{\alpha} \approx \chi_{\Delta} \theta

\]



откуда выводится формула



\[

E_{\theta}(k) \sim \varepsilon_{\theta} \varepsilon^{2 / 3} \chi^{-3} k^{-17 / 3}

\]



В стратифицированных средах К. а. турбулентности усложняется: к определяющим параметрам добавляется параметр плавучести $\alpha g$, где $\alpha-$ коэффициент расширения жидкости (скажем, термического), а $\boldsymbol{g}$ - ускорение силы тяжести. В мелкомасштабной части инерционно-конвективного интервала спектра турбулентности архимедовы силы малы по сравнению с силами инерции и действует обычная К. a. Но если в этот интервал попадает «масштаб плавучести»



\[

L_{*}=(\alpha g)^{-3 / 2} \varepsilon^{5 / 4} \varepsilon_{T}^{-3 / 4},

\]



то, по гипотезе подобия Обухова-Болджиано (1959), статистич. характеристики компонент развитой турбулентности с масштабами из этого интервала будут полностью определяться параметрами $\varepsilon, \varepsilon_{T}$ и $\alpha g$. Это означает, в частности, что константы $C_{1}$ и $C_{1}$ в формулах законов пяти третей для спектров энергии и температурного поля теперь оказываются нек-рыми универсальными функциями от безразмерного волнового числа $L_{*} k$ (стремящимися к вышеуказанным константам на мелкомасштабном конце инерционно-конвективного интервала спектра, то есть асимптотически при возрастании $L_{*} k$ ).



При устойчивой стратификации энергия турбулентности, передаваемая по спектру от болыших масштабов к малым (колмогоровский каскадный процесс), затрачивается в основном на работу против архимедовых сил, и лишь очень





\end{document}